\noindent\textbf{Por qué se usa este procedimiento.} El criterio de Tresca utiliza la mayor diferencia entre las tres principales, incluyendo la nula fuera del plano en tensión plana. Se evalúan todos los pares, se toma su máximo absoluto y se divide el límite uniaxial por ese valor.

\begin{enumerate}

\item La mayor diferencia es $120 - (-40) = 160$ MPa.

\end{enumerate}

\noindent\textbf{Operaciones y resultados intermedios.}

\begin{enumerate}

\item Coeficiente de seguridad.

\[\fitmath{y_{1}=\frac{430}{120--40}\approx 2.6875\,}\]

\end{enumerate}
